% Sinh tự động — scripts/make_report_data.py
\begin{table}[htbp]
  \centering
  \footnotesize
  \setlength{\tabcolsep}{4pt}
  \begin{tabular}{@{}l l l l@{}}
    \toprule
    ID & Hàm mục tiêu & Nguồn ảnh cho xoay & Nguồn ảnh cho loss bổ sung \\
    \midrule
    \cfg{E3} & $\Lsup + \lamc \Lntxent$ & Không có nhánh xoay & $\Dtrain$; không dùng nhãn siêu lớp \\
    \cfg{E4} & $\Lsup + \lamc \Lsupcon$ & Không có nhánh xoay & Chỉ $\DL$ có nhãn hợp lệ \\
    \cfg{E5} & $\Lsup + \lamrot \Lrot + \lamc \Lntxent$ & $\Dtrain$ & NT-Xent: $\Dtrain$; rotation: $\Dtrain$ \\
    \cfg{E6} & $\Lsup + \lamrot \Lrot + \lamc \Lsupcon$ & $\Dtrain$ & SupCon: $\DL$; rotation: $\Dtrain$ \\
    \bottomrule
  \end{tabular}
  \caption{Bốn cấu hình mở rộng contrastive. $\lamc$ dùng chung cho mọi cấu hình và được khóa bằng validation trước khi đánh giá test (\secref{sec:results-ext-lambda}). SupCon chỉ dùng $\DL$; NT-Xent dùng $\Dtrain$ vì không cần nhãn siêu lớp.}
  \label{tab:extensions}
\end{table}